\introduction

\introduction


% 1.Define the task
% What? Ensembles of climate models. Why needed?
For many decades, climate models have been an important tool for gaining a better understanding of the climate system, and for making inferences about future climate and the consequences of climate change. Over the years many different models have become available which is generally beneficial to get a more accurate picture of reality. There are, for instance, different possibilities on how to parameterize processes that cannot be resolved in the models, which leads to uncertainties in the predictions that can be represented by using an ensemble of models. This type of uncertainty is referred to as \emph{model}- or \emph{structural uncertainty}, and it is the primary uncertainty addressed by a multi-model ensemble.
In the most straightforward way, each model in an ensemble is considered equally important and thus receives equal weight. This is therefore known as the ``model democracy'' approach, where simple multi-model means (MMMs) are computed and the ensemble spread is used to quantify uncertainty. While this has long been the standard approach, it comes along with several problems that can be addressed by weighting the models in the ensemble.

% What are the problems with ensembles?
A general challenge is that we must work with ``ensembles of opportunity'': the available models do not represent a statistically perfect ensemble in the sense of a set of unbiased and independent random samples \citep[e.g., see][]{massonClimateModelGenealogy2011, pennellEffectiveNumberClimate2011}. They were not constructed to cover the entire range of plausible outcomes, nor is this space evenly sampled as the models can be dependent on one another in complex and often intransparent ways. Due to such inter-model dependencies, they may, for instance, have shared biases, which essentially means that some models in the ensemble do not provide as much new information as others \citep[e.g.,see][]{junSpatialAnalysisQuantify2012}. Furthermore, the models may differ in how plausible they are in light of the observational data. Therefore, there are two main aspects that weights assigned to models in an ensemble should account for: model performance and model dependence.

% Context: general methods from the literature
Various methods have been proposed in the literature to compute model weights that reflect one or both aspects. 
A method that has received much attention is ClimWIP from \citet{brunnerReducedGlobalWarming2020}, which was developed from a long line of previous work \citep[e.g.,][]{sandersonRepresentativeDemocracyReduce2015, sandersonAddressingInterdependencyMultimodel2015,knuttiClimateModelProjection2017, lorenzProspectsCaveatsWeighting2018}.
It combines weights that are calculated separately for independence, defined in terms of model output similarities, and performance.
Both types of weights are computed based on the assumption of a Gaussian error model and by taking one or more diagnostic variables into account. The diagnostic variables used for computing performance weights may be different from those used to calculate independence weights. 
Meanwhile, the Reliability Ensemble Averaging method (REA) is a similar method introduced by  \citet{giorgiCalculationAverageUncertainty2002a} which also considers model performance and independence.
Another, conceptually different approach is what is often referred to as Bayesian Model Averaging (BMA). While the other methods mentioned above compute model weights based on the performance and/or independence of the individual models in the ensemble, BMA-weights are based on the performance of the weighted ensemble itself. Therefore, the BMA-weights are selected so that the weighted average model simulates the observed data as closely as possible \citep[e,g.][]{massoudBayesianModelAveraging2020,woottenEffectStatisticalDownscaling2020,massoudBayesianWeightingClimate2023,woottenAssessingSensitivitiesClimate2023}. 
It should be noted that this definition of BMA differs from what BMA refers to in the context of statistical models where it was first applied, but since the application of BMA on forecasting surface temperatures by \citet{rafteryUsingBayesianModel2005}, it has become a common way to refer to the method used in at least some climate science studies today.

Although the methods differ in the details, they can mainly be distinguished by whether the ensemble is considered as a whole, which we will refer to as ``\emph{Ensemble Performance Weighting}'', or as individual members, which we will refer to as ``\emph{Individual Performance Weighting}``.
% Questions to be addressed
All approaches to computing model weights have in common that they need to provide answers to the following two questions: (i) how to compute model performance and (ii) how to account for potential dependencies between models.
% Examples for ways to compute performance
Concerning model performance, \citet{sierraInformationtheoreticApproachObtain2025}, for instance, use a method, embedded in information-theoretic concepts, that computes performance weights relative to the inverse of the deviations between model output and data. Similarly, \citet{fengComparisonFourEnsemble2011} use inverse-distance power weighting ($\frac{1}{\textrm{AE}_i^r}, r>0$, AE$_i$: absolute error of model $i$) to compute performance weights. 
Others define performance based on the assumption of independent and identically distributed Gaussian errors \citep[e.g.~][]{brunnerReducedGlobalWarming2020}.
Note that all such individual-performance weighting approaches where the weights are a function of the sum of squared errors yield the same ranking order, only the magnitude of the weights differs (see Appendix \ref{sec:appendix-B} for examples).

% Examples for ways to compute independence
Meanwhile, independence is often defined in terms of model similarities \citep[e.g.~see][]{brunnerReducedGlobalWarming2020}. The more similar the predictions of two models are, the more dependent they are considered to be. This is reasonable in that two dependent models will have shared biases which is reflected in similar model predictions for both. 
Another option is to use meta information about the climate models that could provide further insights about co-dependencies between the models, e.g., information about the modeling group, specific model components or model resolution  \citep[e.g.~see][]{boeInterdependencyMultimodelClimate2018, kumaClimateModelCode2023}. 
However, contrary to the model predictions, this information is not always available, making model similarities a pragmatic choice to define model independence. Furthermore, even if meta information about the models were available, there would not be a unique straightforward definition of independence. As pointed out by \citet{massoudBayesianWeightingClimate2023}, considering full distributions over weight vectors rather than a  single weight per model implicitly provides information about inter-model dependencies. The idea is that this is due to the correlation between the posterior weights of different models. This is an important claim that is fundamental to the distinction between ‘individual’ performance weighting and ‘ensemble’ performance weighting. 

% Short summary paragraph, hinting towards now the main methods sections that come below.
Here we aim to clearly define these two approaches within the common language of probabilities in Bayesian terms, so that we can easily compare and apply both methods in a consistent way. We use simple, idealized examples to test the methods in a systematic way, with a particular focus on the question of how inter-model dependencies are accounted for in a model weighting exercise.  We also  explore how  the results obtained from applying this approach are affected by different methodological choices about the prior distribution over weights, which can be important and should be a prominent aspect of the ensemble performance weighting approach. 

The paper is structured as follows. In Section \ref{sec:framing} we lay out the differences between the two approaches and define the terminology we use. In Sections \ref{sec:priors} and \ref{sec:likelihoods}, we discuss choices for the prior and likelihood distributions and consider the influence of using a set of different diagnostic variables, including a simple example. In Section \ref{sec:dependence} we consider the issue of model dependence through simple examples, while in Section \ref{sec:ecs}, we consider the particular realistic case of using Equilibrium Climate Sensitivity (ECS) as a performance diagnostic. In Section \ref{sec:discussion}, we discuss the overall implications of our work. Section \ref{sec:conclusion} concludes the paper and provides an outlook on future work.




